% Flavor catalog per-process template.
% process_id: E007
% family: edm_neutrino
% owner_agent_id: PKA-E007
% writer_agent_id: null
% checker_agent_id: null
% opus_signoff_id: null
% Canonical metadata sidecar: E007.yaml
% Status is controlled by the YAML sidecar status_history, not this comment.

\section{E007: Radium and xenon electric dipole moments}

\subsection*{Process}
\textbf{Standard notation:} \(d_{\mathrm{Ra}}\) for
\(^{225}\mathrm{Ra}\) and \(d_{\mathrm{Xe}}\) for
\(^{129}\mathrm{Xe}\).  This entry covers atomic EDM searches in radium and
xenon, cataloged in the \texttt{edm\_neutrino} family as complementary
diamagnetic-atom and nuclear CP probes beyond the neutron and
\(^{199}\mathrm{Hg}\) benchmark constraints.

\subsection*{PDG / equivalent value(s)}
No standalone PDG Live Ra/Xe atom-EDM datablock was located for this PKA
deposit, so the direct experiment papers are used as PDG-equivalent sources.
For \(^{225}\mathrm{Ra}\), the canonical Argonne direct limit is in the
sidecar snapshot \texttt{bishof2016\_arxiv1606\_04931.txt},
\[
  |d(^{225}\mathrm{Ra})| < 1.4\times10^{-23}\ e\,\mathrm{cm}
  \qquad (95\%~\mathrm{C.L.}),
\]
from \texttt{flavor\_catalog/references/E007/bishof2016\_arxiv1606\_04931.txt}.
The earlier proof-of-principle Ra measurement in
\texttt{parker2015\_arxiv1504\_07477.txt} gave
\(|d(^{225}\mathrm{Ra})|<5.0\times10^{-22}\ e\,\mathrm{cm}\) at
\(95\%\) confidence and is kept as a superseded post-2008 milestone.

For \(^{129}\mathrm{Xe}\), the strongest direct limit in the deposited
sources is \texttt{sachdeva2019\_arxiv1909\_12800.txt},
\[
  d_A(^{129}\mathrm{Xe}) =
  (1.4 \pm 6.6_{\rm stat} \pm 2.0_{\rm syst})\times10^{-28}\ e\,\mathrm{cm},
  \qquad
  |d_A(^{129}\mathrm{Xe})| < 1.4\times10^{-27}\ e\,\mathrm{cm}
  \quad (95\%~\mathrm{C.L.}).
\]
An independent Heidelberg-centered \(^{129}\mathrm{Xe}\) measurement in
\texttt{allmendinger2019\_arxiv1904\_12295.txt} gives
\((-4.7\pm6.4)\times10^{-28}\ e\,\mathrm{cm}\) and
\(|d_{\mathrm{Xe}}|<1.5\times10^{-27}\ e\,\mathrm{cm}\) at
\(95\%\) C.L.

\subsection*{Relevance to the RS / anarchic-flavor pipeline}
Warped anarchic flavor models generically contain new CP phases.  Loop
matching can generate quark EDMs, chromo-EDMs, the Weinberg three-gluon
operator, and semileptonic CP-odd operators.  Diamagnetic atoms do not map
directly onto one elementary dipole; their dominant hadronic response is
usually expressed through nuclear Schiff moments and CP-odd hadronic
interactions.  \(^{225}\mathrm{Ra}\) is especially useful because octupole
deformation enhances Schiff-moment sensitivity, while \(^{129}\mathrm{Xe}\)
offers a mature noble-gas co-magnetometer platform with different nuclear and
atomic systematics.

\subsection*{Post-2008 developments}
Relative to the 2008 warped-flavor baseline of
\texttt{CsakiFalkowskiWeiler:RSFlavor2008}, Ra and Xe EDMs moved from mostly
future context into direct experimental inputs.  Argonne reported the first
\(^{225}\mathrm{Ra}\) EDM measurement in 2015
\texttt{parker2015\_arxiv1504\_07477.txt} and improved it in 2016 by a
factor of \(36\) \texttt{bishof2016\_arxiv1606\_04931.txt}.  The 2016 Ra
paper also records planned
sensitivity at the \(1\times10^{-28}\ e\,\mathrm{cm}\) statistical level with
\(4\times10^{-29}\ e\,\mathrm{cm}\) systematic uncertainty
\texttt{bishof2016\_arxiv1606\_04931.txt}.  The deposited Argonne program page
\texttt{Argonne:Ra225EDMPage} states the physical motivation:
\(^{225}\mathrm{Ra}\) has a \(15\) day half-life, spin \(1/2\), and expected
sensitivity to T-odd, P-odd effects \(2\)--\(3\) orders of magnitude larger
than \(^{199}\mathrm{Hg}\).

For Xe, the 2019 comagnetometer measurements
\texttt{sachdeva2019\_arxiv1909\_12800.txt} and
\texttt{allmendinger2019\_arxiv1904\_12295.txt} established
\(10^{-27}\ e\,\mathrm{cm}\)-level direct limits.  The DFG
Heidelberg upgrade proposal \texttt{DFG:Xe129Upgrade2026} for 2021--2026
aims to improve the present \(1.5\times10^{-27}\ e\,\mathrm{cm}\) limit by
more than \(2\) orders of magnitude.  KU Leuven's 2023--2027 RaF/AcF
molecular project \texttt{KULeuven:RaFMoleculeProject2023} is not a direct
Ra-atom limit, but it is relevant future context because embedding
octupole-enhanced Ra or Ac nuclei in dipolar molecules can enhance
Schiff-moment sensitivity by several orders of magnitude.

\subsection*{Constraint validity and model dependence}
The direct atom-EDM limits are robust null searches for P- and T-violating
energy shifts.  Their interpretation in an RS pipeline is highly
model-dependent: one must choose a CP-odd low-energy operator basis, run and
match through QCD thresholds, and adopt hadronic, nuclear, and atomic
coefficients.  These entries should therefore be treated as EDM-adjacent
loop/nuclear-translation probes, not as direct \(\Delta F=2\) quark-scan
constraints.

\subsection*{Code coverage in this repo}
\textbf{NO.}  The sidecar \texttt{code\_coverage} block records the required
greps over \texttt{quarkConstraints/}, \texttt{qcd/},
\texttt{flavorConstraints/}, \texttt{neutrinos/}, \texttt{yukawa/},
\texttt{warpConfig/}, \texttt{solvers/}, \texttt{scanParams/}, and
\texttt{tests/} found no Ra, Xe, Schiff-moment, diamagnetic-atom, chromo-EDM,
or Weinberg-operator implementation.  The only nearby dipole code is the
off-diagonal \(\mu\to e\gamma\) helper at
\texttt{flavorConstraints/muToEGamma.py:3},
\texttt{flavorConstraints/muToEGamma.py:21}, and
\texttt{flavorConstraints/muToEGamma.py:81}; that code is not an atomic or
hadronic EDM calculation.

\subsection*{Implementation difficulty}
\textbf{HIGH.}  A live E007 constraint needs new CP-odd quark/gluon and
semileptonic matching, RG evolution to hadronic scales, and separate
nuclear/atomic response layers for Ra and Xe.  The existing
SLL/SLR/VLL/VRR/LR \(\Delta F=2\) basis and the \(\mu\to e\gamma\) dipole
helper do not cover this observable.

\subsection*{Key references}
Process-local source keys before bibliography consolidation:
\texttt{Bishof:Ra225EDM2016}, \texttt{Parker:Ra225EDM2015},
\texttt{Sachdeva:Xe129EDM2019}, \texttt{Allmendinger:Xe129EDM2019},
\texttt{Argonne:Ra225EDMPage}, \texttt{DFG:Xe129Upgrade2026},
\texttt{KULeuven:RaFMoleculeProject2023}, and
\texttt{CsakiFalkowskiWeiler:RSFlavor2008}.
